\textbf{Phát biểu.} Cho \(0<\varepsilon<1\), \(n\ge2\) điểm \(x_1,\dots,x_n\in\mathbb R^d\), số nguyên \(k\ge\dfrac{4\ln n}{\varepsilon^2/2-\varepsilon^3/3}\), ma trận \(R\in\mathbb R^{k\times d}\) có các phần tử \(R_{ij}\sim\mathcal N(0,1)\) độc lập, và \(f(x)=Rx/\sqrt k\). Khi đó với xác suất ít nhất \(1/n\), mọi cặp \(u,v\) trong \(n\) điểm thỏa \[(1-\varepsilon)\|u-v\|^2\le\|f(u)-f(v)\|^2\le(1+\varepsilon)\|u-v\|^2 .\]

\textbf{Chứng minh.} \emph{Bước 1 --- phân phối của một cặp.} Cố định \(u\ne v\), đặt \(z=u-v\). Do \(f\) tuyến tính, \(f(u)-f(v)=Rz/\sqrt k\). Hàng thứ \(j\) của \(R\) là \(r_j\sim\mathcal N(0,I_d)\) nên \(r_j^\top z\sim\mathcal N(0,\|z\|^2)\), và các \(r_j^\top z\) độc lập theo \(j\). Vậy \[Y:=\frac{\|Rz\|^2}{\|z\|^2}=\sum_{j=1}^k\Big(\frac{r_j^\top z}{\|z\|}\Big)^2\sim\chi^2_k,\qquad \frac{\|f(u)-f(v)\|^2}{\|u-v\|^2}=\frac Yk .\] (Nếu \(u=v\) thì bất đẳng thức hiển nhiên đúng.)

\emph{Bước 2 --- hàm sinh mômen.} Với \(g\sim\mathcal N(0,1)\) và \(t<1/2\): \(\mathbb E e^{tg^2}=\int\frac{1}{\sqrt{2\pi}}e^{-(1-2t)g^2/2}dg=(1-2t)^{-1/2}\). Do độc lập, \(\mathbb Ee^{tY}=(1-2t)^{-k/2}\).

\emph{Bước 3 --- đuôi trên.} Với \(0<t<1/2\), theo bất đẳng thức Markov cho \(e^{tY}\): \(\Pr[Y\ge(1+\varepsilon)k]\le e^{-t(1+\varepsilon)k}(1-2t)^{-k/2}\). Chọn \(t=\frac{\varepsilon}{2(1+\varepsilon)}\) (khi đó \(1-2t=\frac1{1+\varepsilon}\)): \[\Pr[Y\ge(1+\varepsilon)k]\le\big((1+\varepsilon)e^{-\varepsilon}\big)^{k/2}=\exp\Big(\tfrac k2\big(\ln(1+\varepsilon)-\varepsilon\big)\Big)\le\exp\Big(-\tfrac k2\big(\tfrac{\varepsilon^2}2-\tfrac{\varepsilon^3}3\big)\Big),\] dùng \(\ln(1+\varepsilon)\le\varepsilon-\frac{\varepsilon^2}2+\frac{\varepsilon^3}3\) với \(\varepsilon\ge0\) (hàm \(h(\varepsilon)=\varepsilon-\frac{\varepsilon^2}2+\frac{\varepsilon^3}3-\ln(1+\varepsilon)\) có \(h(0)=0\), \(h'(\varepsilon)=\frac{\varepsilon^3}{1+\varepsilon}\ge0\)).

\emph{Bước 4 --- đuôi dưới.} Với \(t>0\): \(\Pr[Y\le(1-\varepsilon)k]=\Pr[e^{-tY}\ge e^{-t(1-\varepsilon)k}]\le e^{t(1-\varepsilon)k}(1+2t)^{-k/2}\). Chọn \(t=\frac{\varepsilon}{2(1-\varepsilon)}\) (\(1+2t=\frac1{1-\varepsilon}\)): \[\Pr[Y\le(1-\varepsilon)k]\le\big((1-\varepsilon)e^{\varepsilon}\big)^{k/2}\le\exp\Big(-\tfrac{k\varepsilon^2}4\Big)\le\exp\Big(-\tfrac k2\big(\tfrac{\varepsilon^2}2-\tfrac{\varepsilon^3}3\big)\Big),\] dùng \(\ln(1-\varepsilon)\le-\varepsilon-\frac{\varepsilon^2}2\) (khai triển Taylor, mọi số hạng âm) và \(\varepsilon^3/3\ge0\).

\emph{Bước 5 --- chặn hợp.} Với \(k\ge\frac{4\ln n}{\varepsilon^2/2-\varepsilon^3/3}\) (mẫu số dương vì \(\varepsilon<1\)): \(\frac k2(\frac{\varepsilon^2}2-\frac{\varepsilon^3}3)\ge2\ln n\), nên mỗi đuôi \(\le n^{-2}\) và một cặp bị méo với xác suất \(\le 2/n^2\). Có \(\binom n2\) cặp, theo chặn hợp: \[\Pr[\exists\text{ cặp bị méo}]\le\binom n2\frac2{n^2}=\frac{n-1}n=1-\frac1n .\] Vậy xác suất mọi cặp đều thỏa ít nhất \(1/n\). \(\blacksquare\)

\emph{Nhận xét.} Xác suất \(1/n\) có vẻ nhỏ, nhưng lặp lại \(O(n)\) phép chiếu độc lập là gần như chắc chắn thành công; trên thực tế chặn này rất bảo thủ (xem C.3). Điểm mấu chốt: \(k\) \textbf{chỉ phụ thuộc \(\ln n\) và \(\varepsilon\), không phụ thuộc \(d\)}.
